\documentclass[a4paper,10pt,twoside,openany]{article}

\usepackage[lang=hebrew]{maths}
\usepackage{hebrewdoc}
\usepackage{stylish}
\usepackage{lipsum}
\let\bs\blacksquare

\usepackage{siunitx,array}
\sisetup{table-format=-1.5, group-digits=false}
\newcolumntype{L}{>{\phantom{$-$}}l}       % prefix some whitespace
\newcommand\mcL[1]{\multicolumn{1}{L}{#1}} % handy shortcut macro

\setlength{\parindent}{0pt}

%%%%%%%%%%%%
% Styling %
%%%%%%%%%%%%

\usepackage{enumitem}

%%%%%%%%%%%%%
% Counters  %
%%%%%%%%%%%%%

\setcounter{section}{1}     
            
%BIBLIOGRAPHY
\usepackage[
backend=biber,
style=alphabetic,
]{biblatex}
\addbibresource{bibliography.bib} %Imports bibliography file

%%%%%%%%%%
% Title  %
%%%%%%%%%%
\title{
אלגברה ב' - גיליון תרגילי בית 10 \\
עוד אופרטורים על מרחבי מכפלה פנימית, ומשפט הפירוק הספקטרלי
\\
\vspace{1cm}
\large{תאריך הגשה: 9.7.2026}
}
\date{}

\begin{document}
\maketitle

\textbf{לאורך הגיליון, נניח כי
$V$
מרחב וקטורי סוף־מימדי מעל שדה
$\FF$.}

\begin{notation}
עבור
$T \in \End_\FF\prs{V}$
נסמן ב־%
$\sigma\prs{T}$
את קבוצת הערכים העצמיים של
$T$
ונסמן
\begin{align*}
\text{.} r_\sigma\prs{T} \coloneqq \max_{\lambda \in \sigma\prs{T}} \abs{\lambda}
\end{align*}
\end{notation}

\begin{exercise}
יהי
$T \in \End_\CC\prs{V}$
נורמלי
ויהיו
$p,q \in \CC\brs{x}$
כך שהאופרטור
$q\prs{T}$
הפיך.
נסמן
$S = p\prs{T} q\prs{T}^{-1}$.

הוכיחו כי
\begin{align*}
\sigma\prs{S} = \set{\frac{p\prs{\lambda}}{q\prs{\lambda}}}{\lambda \in \sigma\prs{T}}
\end{align*}
וכי לכל
$\lambda' \in \sigma\prs{S}$
מתקיים
\begin{align*}
r^S_a\prs{\lambda'} = \sum_{\substack{\lambda \in \sigma\prs{T} \\ \frac{p\prs{\lambda}}{q\prs{\lambda}} = \lambda'}} r_a^T\prs{\lambda} \\
\text{.} r^S_g\prs{\lambda'} = \sum_{\substack{\lambda \in \sigma\prs{T} \\ \frac{p\prs{\lambda}}{q\prs{\lambda}} = \lambda'}} r_g^T\prs{\lambda}
\end{align*}
\end{exercise}

\begin{exercise}
יהי
$T \in \End_{\mbb{C}}\prs{V}$
נורמלי, ונניח כי
$3,4$
ערכים עצמיים של
$T$.
הראו שיש
$v \in V$
עבורו
$\norm{v} = \sqrt{2}$
וגם
$\norm{Tv} = 5$.
\end{exercise}

\begin{exercise}
יהי
$T \in \End_{\mbb{C}}\prs{V}$
ויהי
$a \in \mbb{C}$
עבורו
$\abs{a} \neq 1$.

\begin{enumerate}
\item הראו כי אם
$T^* = aT$
אז
$T = 0$.

\item נניח כי
$T$
נורמלי, ויהי
$S = T - a T^*$.
הוכיחו כי
$\ker\prs{T} = \ker\prs{S}$.
\end{enumerate}
\end{exercise}

\begin{exercise}
\begin{enumerate}
\item
יהי
$T \in \End_\CC\prs{V}$
הרמיטי עם
$r_\sigma\prs{T} \leq 1$.
מיצאו
$S \in \End_\CC\prs{V}$
אוניטרי עבורו
\begin{align*}
\text{.} T = \frac{S + S^*}{2}
\end{align*}

\item יהי
$T \in \End_\CC\prs{V}$.
הוכיחו כי
$T$
ניתן לכתיבה כצירוף לינארי של ארבעה אופרטורים אוניטריים.
\end{enumerate}
\end{exercise}

\begin{exercise}
\textbf{שיקוף}
ב־%
$\RR^n$
הוא אופרטור
$T$
כך שקיים
$v \in V$
עבורו
$T\prs{v} = -v$
וגם
$T\prs{w} = w$
לכל
$w \in \Span\prs{v}^\perp$.
נסמנו
$R_v$
ונסמן גם
$R \coloneqq R_{e_1}$.

\begin{enumerate}
\item תהי
$\theta \in \RR$,
תהי
\begin{align*}
A_\theta \coloneqq \pmat{\cos\prs{\theta} & -\sin\prs{\theta} \\ \sin\prs{\theta} & \cos\prs{\theta}}
\end{align*}
ויהי
$\rho_\theta \coloneqq T_{A_\theta}$
אופרטור הכפל משמאל ב־%
$A_\theta$.

מיצאו
$v \in V$
עבורו
$\rho_\theta = R_v \circ R$
והסיקו כי כל אופרטור אורתוגונלי על
$\RR^2$
הוא הרכבה של לכל היותר שני שיקופים.

\item
יהי
$T \in \End_\RR\prs{\RR^3}$
אורתוגונלי. הראו שקיים בסיס
$\mcal{B}$
של
$\RR^3$,
שקיימת
$\lambda \in \set{\pm 1}$
ושקיימת
$O \in \Mat_2\prs{\RR}$
אורתוגונלית עבורם מתקיים
\begin{align*}
\text{.} \brs{T}_{\mcal{B}} = \pmat{\lambda & 0_{1 \times 2} \\ 0_{2 \times 1} & O}
\end{align*}
\item הוכיחו כי כל אופרטור אורתוגונלי
$T \in \End_\RR\prs{\RR^3}$
הוא הרכבה של לכל היותר
$3$
שיקופים.
\end{enumerate}
\end{exercise}

\end{document}